\documentclass{article}

\usepackage{float}
\restylefloat{table}

\usepackage{booktabs}

\title{Team Productivity: POC\\\progname}

\author{\authname}

\date{}

\input{../Comments.text}
\input{../Common.text}

\begin{document}

\maketitle

This document summarizes the contributions of each team member up to the POC
Demo.  The time period of interest is the time between the beginning of the term
and the POC demo.


\section{Demo Plans}

The main feasibility risk is whether the network can produce an accurate aggregate
and update provider reputations when ground truth is unavailable to the protocol.
In particular, biased or coordinated providers may gain influence by agreeing with
one another, especially when the network starts without established reputations.

We will demonstrate an off-chain simulation of repeated linear-regression inference
requests. Honest providers submit noisy predictions, biased providers submit
inaccurate predictions, and learners improve their predictions over time. We will
also examine providers that share a bias or coordinate their predictions. In each
round, providers submit predictions, the protocol aggregates them, and their
reputations are updated. Ground truth is used afterward to evaluate the outcome;
it is not used for aggregation or reputation updates.

The demo will examine the following scenarios, under the assumption that honest providers form a
majority and fewer than half of the providers behave maliciously:
\begin{itemize}
    \item With honest and biased providers starting at equal reputations, honest
    providers should gain reputation on average, biased providers should lose
    reputation on average, and average relative error between the aggregate and
    ground truth should decrease over successive rounds.
    \item When learners improve their predictions, their reputations should rise
    after the improvement, consistently inaccurate providers should lose
    reputation, and average relative error should decrease.
    \item When a minority of providers submit similar inaccurate predictions or
    coordinate their submissions, their combined reputation should decline
    rather than increase through agreement, while average relative error should
    decrease over time.
\end{itemize}

For each scenario, we will show provider reputations and relative
error across rounds of the simulation. We will also compare the reputation-weighted
aggregate with an unweighted median to assess the effect of reputation on the aggregate.

\section{Team Meeting Attendance}

\begin{table}[H]
\centering
\begin{tabular}{ll}
\toprule
\textbf{Student} & \textbf{Meetings}\\
\midrule
Total & 9\\
Robert Lam & 9\\
Nathan Hum & 9\\
Aidan Lao & 9\\
Keira Laskoski & 9\\
\bottomrule
\end{tabular}
\end{table}

Counts are taken from the team meeting issues in the repository (\#1, \#41,
\#58, \#59, \#81, \#107, \#114, \#185 and \#214).

\section{Supervisor/Stakeholder Meeting Attendance}

\noindent \textbf{Supervisor's Name: } Dr. Mohamed Hassan

\begin{table}[H]
\centering
\begin{tabular}{ll}
\toprule
\textbf{Student} & \textbf{Meetings}\\
\midrule
Total & 3\\
Robert Lam & 3\\
Nathan Hum & 3\\
Aidan Lao & 2\\
Keira Laskoski & 3\\
\bottomrule
\end{tabular}
\end{table}

Counts are taken from the supervisor meeting issues in the repository (\#71,
\#94 and \#146). Aidan did not attend the meeting in issue \#94, but this expected absence was 
known at the time of making the meeting. Due to the supervisor's schedule, there was no
other time that the meeting could be scheduled.

\section{Lecture Attendance}

\begin{table}[H]
\centering
\begin{tabular}{ll}
\toprule
\textbf{Student} & \textbf{Lectures}\\
\midrule
Total & 11\\
Robert Lam & 5 \\
Nathan Hum & 5 \\
Aidan Lao & 5 \\
Keira Laskoski & 5 \\
\bottomrule
\end{tabular}
\end{table}

All team members attended the first five lectures at the beginning of the term. 
Counts are taken from lecture issues in the repository (\#254, \#253, \#252, \#250, and \#251).

The remaining lectures were not attended. However, team members reviewed the lecture 
slides individually and discussed them as needed during team meetings.

\section{TA Document Discussion Attendance}

\noindent \textbf{TA's Name: } Rashad Bhuiyan

\begin{table}[H]
\centering
\begin{tabular}{ll}
\toprule
\textbf{Student} & \textbf{Meetings}\\
\midrule
Total & 4\\
Robert Lam & 4\\
Nathan Hum & 4\\
Aidan Lao & 4\\
Keira Laskoski & 4\\
\bottomrule
\end{tabular}
\end{table}

Counts are taken from the TA meeting issues in the repository (\#39, \#80,
\#147 and \#215). Every team member attended every meeting.

\section{Commits, Lines Added and Lines Deleted, Pull Requests}

\begin{table}[H]
\centering
\begin{tabular}{lllllll}
\toprule
\textbf{Student} & \textbf{Commits} & \textbf{Percent} & \textbf{Lines Added} &
\textbf{Lines Deleted} & \textbf{Pull Requests}\\
\midrule
Total & 189 & 100\% & 52338 & 26463 & 120\\
Robert Lam & 62 & 32.8\% & 19061 & 5814 & 34\\
Nathan Hum & 31 & 16.4\% & 16044 & 13571 & 32\\
Aidan Lao & 35 & 18.5\% & 9373 & 2027 & 24\\
Keira Laskoski & 61 & 32.3\% & 7860 & 5051 & 29\\
\bottomrule
\end{tabular}
\end{table}

Commit counts are non-merge commits reachable from \texttt{main}, combined
across each member's Git identities. The 110 commits made by the
\texttt{github-actions[bot]} account to publish generated PDFs are excluded.

Line counts include documentation as well as code and are inflated by a few
whole-file rewrites. Nathan's commit \texttt{12d39ab} changed the line endings
of 97 files (about 11{,}350 lines added and 11{,}330 deleted), and Robert's
initial commit added the 11{,}354-line course template. Keira's commit
\texttt{1c9f4c5} restored the line endings of \texttt{SRS.tex} (2{,}123 lines
added and 2{,}123 deleted).

\section{Issue Tracker}

\begin{table}[H]
\centering
\begin{tabular}{lll}
\toprule
\textbf{Student} & \textbf{Authored (O+C)} & \textbf{Assigned (C only)}\\
\midrule
Robert Lam & 29 & 28 \\
Nathan Hum & 5 & 10 \\
Aidan Lao & 5 & 17 \\
Keira Laskoski & 48 & 14 \\
\bottomrule
\end{tabular}
\end{table}

The repository has 122 issues in total. Members of Teams 4 and 23 authored the
remaining 35 issues as peer reviews. Keira's authored count is high because she
created the per-section issues for the Problem Statement, Development Plan and
Team Charter, and logged the TA feedback as issues.

\section{CICD via GitHub Actions}

The project currently has one GitHub Actions workflow that builds the pdf documents from the latex files. 

As described in the Development Plan and VnV Plan, the CI pipeline will be
extended to cover the code in \texttt{src/}. On each pull request, it will build
the Rust workspace, run the unit and system tests with \texttt{cargo test},
check formatting with \texttt{rustfmt}, and lint with \texttt{clippy}. It will
also build and test the Solana programs (using Anchor) and the React frontend,
and report test coverage. Together with the requirement that every pull request
is approved by at least one other team member before it is merged, this will
stop changes that break the build, the tests or the coding standard from
reaching \texttt{main}. 

\section{Custom Documents (CDs)}

The Fall CD was the Cognitive Walkthrough, which has been completed.

For the winter term, the current plan is to write a Business Plan. The team believes this will
benefit the project, since token economics and economical incentives are deeply integrated into the 
foundation of the DAIIM project.

However, the team may reasses this plan in the following few weeks.


\section{Team Charter Trigger Items}

\textbf{Trigger Summary:}
\begin{itemize}
    \item Meeting attendance falling below 95\%.
    \item Issues created without labels for type, milestone, assignee(s), and target deadlines.
    \item Primary reviewer fails to review PR within 48h.
    \item Workload is not distributed evenly.
\end{itemize}

\textbf{Observed Trigger Events:}
\begin{itemize}
    \item There was one instance of a member missing a meeting, but they stated that they would be unavailable for that time before the meeting was scheduled. Some members have been late to meetings, but communicated to the team why they were running late and when they were going to arrive.
    \item Many members have failed to created proper issues for their parts and commits.
    \item Some members have missed soft internal deadlines, without giving much notice ahead of time.
    \item There have been instances of PRs not being reviewed within 48h.
\end{itemize}

\textbf{Plan and Reflection:}
The team has not had many issues with absences or lateness. Team members are efficient at communicating
and arriving as soon as possible, and all excuses for absence/lateness have been reasonable.

It has been recognized that the team could be making better use of the Kanban board and the issue
tracker to better organize deadlines and team contributions. Moving forward, all team members will
plan to make correctly labelled issues related to their parts, and close them with commits. The team 
also plans to be more diligent with assigning specific reviewers to every Pull Request, to clear up 
confusion about who shold be reviewing better implement the 48h rule.

Team members who have missed internal deadlines are aware of their errors and plan to be more diligent 
with meeting future imposed deadlines.

The team also recognizes that there have been no team building activities. This is something that 
will be worked on in the winter term.


\section{Additional Productivity Metrics}

There are no additional metrics to include.

In the upcoming continution of the project, the team plans to count code reviews done and test cases written.


\end{document}